%% kapa: exam=1 round=33 year=2022 period=1 form=A subject=civil_law subject_ko=민법 seq=16 number=16
%% source: Q-Net 감정평가사 기출문제 (cst00309) artlSeq=5209997 file=raw/1008/5209997/2022년도 제33회 감정평가사 1차시험 1교시(민법,경제학,부동산학).pdf
%% conversion: kordoc 4.24.51 -> tools/gen1.py (auto); review=proofread-sonnet; flags=none
\begin{question}{16}
법률행위의 무효에 관한 설명으로 옳지 않은 것은? (다툼이 있으면 판례에 따름)
\begin{choices}
  \item 매매계약이 약정된 매매대금의 과다로 인하여 불공정한 법률행위에 해당하는 경우, 무효행위의 전환에 관한 민법 제138조가 적용될 수 있다.
  \item 취소할 수 있는 법률행위를 취소한 후에도 무효인 법률행위의 추인의 요건과 효력으로서 추인할 수 있다.
  \item 법률행위의 일부무효에 관한 민법 제137조는 임의규정이다.
  \item 집합채권의 양도가 양도금지특약을 위반하여 무효인 경우, 채무자는 집합채권의 일부 개별 채권을 특정하여 추인할 수 없다.
  \item 무효인 가등기를 유효한 등기로 전용하기로 한 약정은 특별한 사정이 없는 한 그때부터 유효하고 이로써 그 가등기가 소급하여 유효한 등기로 전환될 수 없다.
\end{choices}
\end{question}
\answer{4}
